\documentclass[11pt,a4paper]{article}
\usepackage[margin=2.0cm]{geometry}
\usepackage{graphicx}
\usepackage{caption}
\usepackage{subcaption}
\usepackage{hyperref}
\usepackage{grffile}
\usepackage{float}
\usepackage{amsmath}
\usepackage{amssymb}
\usepackage{booktabs}
\usepackage{xcolor}
\usepackage{xspace}

% Shortcuts
\newcommand{\pT}{\ensuremath{p_{\mathrm{T}}}\xspace}
\newcommand{\abseta}{\ensuremath{|\eta|}\xspace}
\newcommand{\Zllg}{\ensuremath{Z \to \ell\ell\gamma}\xspace}
\newcommand{\Zeeg}{\ensuremath{Z \to ee\gamma}\xspace}
\newcommand{\Zmumug}{\ensuremath{Z \to \mu\mu\gamma}\xspace}
\newcommand{\sqrts}{\ensuremath{\sqrt{s} = 13.6~\text{TeV}}\xspace}
\newcommand{\fbinv}{\ensuremath{\text{fb}^{-1}}\xspace}
\newcommand{\pbinv}{\ensuremath{\text{pb}^{-1}}\xspace}

\graphicspath{{../output/kinematics/}}

\title{Kinematic Distributions in the $\Zllg$ Photon Shower Shape Study:\\
  MC Normalisation, Data--MC Agreement, and Conversion Rate Discrepancies}
\author{Marta Fernandes}
\date{\today}

\begin{document}
\maketitle

\begin{abstract}
This note documents the investigation of photon kinematic distributions
($\pT$ and $\abseta$) in $\Zllg$ events selected from 108~\fbinv of
ATLAS data collected at $\sqrts$ in 2024, compared to Sherpa MC23e
Monte Carlo simulation.  A significant normalisation discrepancy was traced
to the use of DAOD-level sum-of-weights in the NTupleMaker framework,
which misses events removed by derivation skimming.  Replacing these
with EVNT-level values from AMI resolved the over-estimation.  The
remaining ${\sim}20\%$ normalisation deficit is attributed to missing
$Z{+}\text{jets}$ background and incomplete MC sample coverage.  Shape
comparisons between data and MC reveal good agreement in $\pT$, while
$\abseta$ distributions show conversion-status--dependent discrepancies
linked to a known TRT barrel tracker issue in Run~3.
\end{abstract}

\tableofcontents
\clearpage

% ====================================================================
\section{Introduction}
\label{sec:introduction}
% ====================================================================

The photon shower shape reweighting analysis requires precise data--MC comparisons
of photon properties in $\Zllg$ events.  As a prerequisite, the kinematic
distributions --- photon transverse momentum $\pT$ and pseudorapidity $\abseta$ ---
must demonstrate adequate agreement between data and simulation, both in
shape and in overall normalisation.

This study was motivated by the observation that the MC prediction
significantly exceeded the data yield.  The investigation revealed a
normalisation bug in the MC event weighting, related to the handling of
the sum-of-weights in the presence of DAOD-level skimming.  After correcting
this, residual differences were found in the $\abseta$ distributions that
depend on the photon conversion status, pointing to a known issue with the
TRT barrel tracker in Run~3.

% ====================================================================
\section{Samples and Event Selection}
\label{sec:samples}
% ====================================================================

\subsection{Data}

The analysis uses 108~\fbinv of proton--proton collision data collected by the
ATLAS detector in 2024 at $\sqrts$.  Events are selected from two DAOD
derivations:
\begin{itemize}
  \item \textbf{EGAM3}: $\Zeeg$ candidates,
  \item \textbf{EGAM4}: $\Zmumug$ candidates.
\end{itemize}
The two channels are analysed separately and also combined into an inclusive
$\Zllg$ sample.

\subsection{Monte Carlo}

Simulated $\Zllg$ events are generated with Sherpa~2.2.14 using the MC23e
campaign.  Two samples are used:

\begin{table}[H]
\centering
\caption{MC samples and normalisation parameters.}
\label{tab:mc-samples}
\begin{tabular}{lcccc}
\toprule
DSID & Process & $\sigma$ [pb] & AMI EVNT sumW & Total events \\
\midrule
700770 & Sh\_2214\_eegamma   & 102.60 & $5.269 \times 10^{14}$ & 185\,402\,000 \\
700771 & Sh\_2214\_mumugamma & 102.59 & $5.421 \times 10^{14}$ & 189\,985\,000 \\
\bottomrule
\end{tabular}
\end{table}

Cross-sections are obtained from the PMG cross-section database (\texttt{PMGxsecDB\_mc23.txt}).
The generator filter efficiency is 1.0 for both samples.

\subsection{Event Selection}

The baseline selection follows Francisco's configuration for the cell-energy
reweighting study:

\begin{table}[H]
\centering
\caption{Event selection criteria.}
\label{tab:selection}
\begin{tabular}{lc}
\toprule
Variable & Requirement \\
\midrule
Photon $\pT$ & $> 10$~GeV \\
Photon $|\eta|$ & $< 2.37$, excluding crack $[1.37, 1.52)$ \\
$m_{\ell\ell}$ & $[40, 83]$~GeV \\
$m_{\ell\ell\gamma}$ & $[80, 100]$~GeV \\
$\Delta R(\gamma, \ell)$ & $> 0.4$ (both leptons) \\
\bottomrule
\end{tabular}
\end{table}

Several scenarios are studied, varying the photon conversion status and
isolation requirement:

\begin{table}[H]
\centering
\caption{Selection scenarios.}
\label{tab:scenarios}
\begin{tabular}{lccc}
\toprule
Scenario & Conversion & Photon ID & Isolation \\
\midrule
Baseline (unconverted) & Unconverted only & No ID & Loose \\
Converted              & Converted only   & No ID & Loose \\
All conversions        & Both             & No ID & Loose \\
Tight isolation        & Unconverted only & No ID & Tight \\
\bottomrule
\end{tabular}
\end{table}

% ====================================================================
\section{MC Event Weight and the Normalisation Problem}
\label{sec:normalisation}
% ====================================================================

\subsection{MC Weight Formula}

Each MC event is weighted by:
\begin{equation}
w = w_{\mathrm{mc}} \times w_{\mu} \times \frac{\sigma \times \mathcal{L}}{\sum w_{\mathrm{EVNT}}} \times \mathrm{SF}_{\ell_1} \times \mathrm{SF}_{\ell_2} \times \mathrm{SF}_{\gamma,\mathrm{iso}}
\label{eq:mc-weight}
\end{equation}
where $w_{\mathrm{mc}}$ is the generator weight, $w_{\mu}$ the pileup
reweighting factor, $\sigma$ the cross-section, $\mathcal{L} = 108$~\fbinv
the integrated luminosity, $\sum w_{\mathrm{EVNT}}$ the EVNT-level
sum of generator weights, and the SF terms are lepton and photon isolation
scale factors.

An alternative weight without scale factors is also computed for comparison:
\begin{equation}
w_{\mathrm{no\,SF}} = w_{\mathrm{mc}} \times w_{\mu} \times \frac{\sigma \times \mathcal{L}}{\sum w_{\mathrm{EVNT}}}
\end{equation}

\subsection{The Bug: DAOD-level Sum-of-Weights}
\label{sec:bug}

The NTupleMaker framework, developed to process input AOD files (which are
unskimmed), computes a sum of generator weights and stores it in a histogram
(\texttt{h\_sumW}) in the output ROOT file.  The sum-of-weights is retrieved
from bin~2 of this histogram.

However, the production ntuples used in this analysis were derived from
\textbf{DAOD} (Derived Analysis Object Data) files, which are \textbf{skimmed}:
events that fail the derivation-level selection are removed before reaching
the NTupleMaker.  Consequently, the recorded sum-of-weights only accounts
for events that survived the DAOD skimming --- approximately \textbf{36\%}
of the full EVNT-level sample.

Using this DAOD-level $\sum w$ in the denominator of Equation~\ref{eq:mc-weight}
inflates the per-event weight by a factor of:
\begin{equation}
\frac{\sum w_{\mathrm{EVNT}}}{\sum w_{\mathrm{DAOD}}} \approx 2.77
\end{equation}

This resulted in the MC prediction overshooting the data by a factor of
${\sim}2.2$ (MC/Data $\approx 2.19$), which was immediately visible in the
event-count kinematic distributions.

\subsection{The Fix: AMI EVNT-level Sum-of-Weights}

The fix replaces the dynamically computed DAOD-level sum-of-weights with
the correct EVNT-level values obtained from AMI (the ATLAS Metadata
Interface).  The corrected values are hardcoded in the analysis configuration:

\begin{itemize}
  \item DSID 700770: $\sum w = 5.269047 \times 10^{14}$
  \item DSID 700771: $\sum w = 5.421893 \times 10^{14}$
\end{itemize}

After this correction, the MC/Data ratio drops to ${\sim}0.79$, i.e.\ the
MC now \emph{under}predicts the data by about 20\%.

\subsection{Residual Normalisation Deficit}
\label{sec:residual}

The remaining ${\sim}20\%$ deficit has several contributing causes:
\begin{enumerate}
  \item \textbf{Missing $Z{+}\text{jets}$ background}:
    The MC samples contain only the signal process $Z \to \ell\ell\gamma$.
    The data inevitably includes $Z{+}\text{jets}$ events where a jet is
    misidentified as a photon.  This background is not subtracted and
    contributes additional events in data.
  \item \textbf{Incomplete MC sample coverage}:
    The available MC samples cover 97.59\% of the total generated events;
    some sub-jobs may have failed or not been included.
  \item \textbf{Filter efficiencies}:
    Generator-level filter efficiencies and potential differences in
    the DAOD derivation acceptance are not fully accounted for.
\end{enumerate}

Crucially, this residual normalisation offset does \textbf{not} affect the
shower shape reweighting, which uses \emph{area-normalised} distributions.
The key requirement is that the \emph{shape} of the kinematic distributions
agrees between data and MC.

% ====================================================================
\section{Results: Loose Isolation}
\label{sec:results-loose}
% ====================================================================

Each figure shows the photon $\pT$ and $\abseta$ distributions for data (black
markers), MC with scale factors (red histogram), and MC without scale factors
(blue dashed histogram).  The left column shows absolute event counts
and the right column shows area-normalised distributions.  A ratio
panel (MC/Data) is included below each comparison.

\subsection{Unconverted Photons}

\begin{figure}[H]
\centering
\begin{subfigure}[b]{0.48\textwidth}
  \includegraphics[width=\textwidth]{loose/llgamma/baseline/plots/kinematics_pt.pdf}
  \caption{$\pT$ (events)}
\end{subfigure}\hfill
\begin{subfigure}[b]{0.48\textwidth}
  \includegraphics[width=\textwidth]{loose/llgamma/baseline/plots/kinematics_pt_norm.pdf}
  \caption{$\pT$ (normalised)}
\end{subfigure}\\[1em]
\begin{subfigure}[b]{0.48\textwidth}
  \includegraphics[width=\textwidth]{loose/llgamma/baseline/plots/kinematics_eta.pdf}
  \caption{$\abseta$ (events)}
\end{subfigure}\hfill
\begin{subfigure}[b]{0.48\textwidth}
  \includegraphics[width=\textwidth]{loose/llgamma/baseline/plots/kinematics_eta_norm.pdf}
  \caption{$\abseta$ (normalised)}
\end{subfigure}
\caption{Unconverted photons, loose isolation.
  The $\pT$ shapes agree well between data and MC.  The $\abseta$ distribution
  shows a deficit in MC in the barrel region, particularly at low $|\eta|$.}
\label{fig:loose-uc}
\end{figure}
\clearpage

\subsection{Converted Photons}

\begin{figure}[H]
\centering
\begin{subfigure}[b]{0.48\textwidth}
  \includegraphics[width=\textwidth]{loose/llgamma/converted/plots/kinematics_pt.pdf}
  \caption{$\pT$ (events)}
\end{subfigure}\hfill
\begin{subfigure}[b]{0.48\textwidth}
  \includegraphics[width=\textwidth]{loose/llgamma/converted/plots/kinematics_pt_norm.pdf}
  \caption{$\pT$ (normalised)}
\end{subfigure}\\[1em]
\begin{subfigure}[b]{0.48\textwidth}
  \includegraphics[width=\textwidth]{loose/llgamma/converted/plots/kinematics_eta.pdf}
  \caption{$\abseta$ (events)}
\end{subfigure}\hfill
\begin{subfigure}[b]{0.48\textwidth}
  \includegraphics[width=\textwidth]{loose/llgamma/converted/plots/kinematics_eta_norm.pdf}
  \caption{$\abseta$ (normalised)}
\end{subfigure}
\caption{Converted photons, loose isolation.
  Notably better data--MC agreement in $\abseta$ compared to the
  unconverted case, as discussed in Section~\ref{sec:conversion}.}
\label{fig:loose-conv}
\end{figure}
\clearpage

\subsection{All Conversions}

\begin{figure}[H]
\centering
\begin{subfigure}[b]{0.48\textwidth}
  \includegraphics[width=\textwidth]{loose/llgamma/all_conv/plots/kinematics_pt.pdf}
  \caption{$\pT$ (events)}
\end{subfigure}\hfill
\begin{subfigure}[b]{0.48\textwidth}
  \includegraphics[width=\textwidth]{loose/llgamma/all_conv/plots/kinematics_pt_norm.pdf}
  \caption{$\pT$ (normalised)}
\end{subfigure}\\[1em]
\begin{subfigure}[b]{0.48\textwidth}
  \includegraphics[width=\textwidth]{loose/llgamma/all_conv/plots/kinematics_eta.pdf}
  \caption{$\abseta$ (events)}
\end{subfigure}\hfill
\begin{subfigure}[b]{0.48\textwidth}
  \includegraphics[width=\textwidth]{loose/llgamma/all_conv/plots/kinematics_eta_norm.pdf}
  \caption{$\abseta$ (normalised)}
\end{subfigure}
\caption{All photons (unconverted + converted), loose isolation.
  The inclusive $\abseta$ shape shows intermediate agreement, mixing the two
  sub-populations.}
\label{fig:loose-all}
\end{figure}
\clearpage

% ====================================================================
\section{Results: Tight Isolation}
\label{sec:results-tight}
% ====================================================================

Applying tight isolation reduces the contribution of $Z{+}\text{jets}$
background, where jets faking photons tend to be less isolated.
As suggested by the supervisors, this was investigated to see if
tighter selection improves the data--MC agreement.

\subsection{Unconverted Photons}

\begin{figure}[H]
\centering
\begin{subfigure}[b]{0.48\textwidth}
  \includegraphics[width=\textwidth]{iso_tight/llgamma/baseline/plots/kinematics_pt.pdf}
  \caption{$\pT$ (events)}
\end{subfigure}\hfill
\begin{subfigure}[b]{0.48\textwidth}
  \includegraphics[width=\textwidth]{iso_tight/llgamma/baseline/plots/kinematics_pt_norm.pdf}
  \caption{$\pT$ (normalised)}
\end{subfigure}\\[1em]
\begin{subfigure}[b]{0.48\textwidth}
  \includegraphics[width=\textwidth]{iso_tight/llgamma/baseline/plots/kinematics_eta.pdf}
  \caption{$\abseta$ (events)}
\end{subfigure}\hfill
\begin{subfigure}[b]{0.48\textwidth}
  \includegraphics[width=\textwidth]{iso_tight/llgamma/baseline/plots/kinematics_eta_norm.pdf}
  \caption{$\abseta$ (normalised)}
\end{subfigure}
\caption{Unconverted photons, tight isolation.
  The normalisation gap narrows slightly compared to the loose isolation
  case (Figure~\ref{fig:loose-uc}), consistent with the rejection of
  $Z{+}\text{jets}$ background in data.}
\label{fig:tight-uc}
\end{figure}
\clearpage

\subsection{Converted Photons}

\begin{figure}[H]
\centering
\begin{subfigure}[b]{0.48\textwidth}
  \includegraphics[width=\textwidth]{iso_tight/llgamma/converted/plots/kinematics_pt.pdf}
  \caption{$\pT$ (events)}
\end{subfigure}\hfill
\begin{subfigure}[b]{0.48\textwidth}
  \includegraphics[width=\textwidth]{iso_tight/llgamma/converted/plots/kinematics_pt_norm.pdf}
  \caption{$\pT$ (normalised)}
\end{subfigure}\\[1em]
\begin{subfigure}[b]{0.48\textwidth}
  \includegraphics[width=\textwidth]{iso_tight/llgamma/converted/plots/kinematics_eta.pdf}
  \caption{$\abseta$ (events)}
\end{subfigure}\hfill
\begin{subfigure}[b]{0.48\textwidth}
  \includegraphics[width=\textwidth]{iso_tight/llgamma/converted/plots/kinematics_eta_norm.pdf}
  \caption{$\abseta$ (normalised)}
\end{subfigure}
\caption{Converted photons, tight isolation.
  The agreement remains comparably good as in the loose case
  (Figure~\ref{fig:loose-conv}).}
\label{fig:tight-conv}
\end{figure}
\clearpage

\subsection{All Conversions}

\begin{figure}[H]
\centering
\begin{subfigure}[b]{0.48\textwidth}
  \includegraphics[width=\textwidth]{iso_tight/llgamma/all_conv/plots/kinematics_pt.pdf}
  \caption{$\pT$ (events)}
\end{subfigure}\hfill
\begin{subfigure}[b]{0.48\textwidth}
  \includegraphics[width=\textwidth]{iso_tight/llgamma/all_conv/plots/kinematics_pt_norm.pdf}
  \caption{$\pT$ (normalised)}
\end{subfigure}\\[1em]
\begin{subfigure}[b]{0.48\textwidth}
  \includegraphics[width=\textwidth]{iso_tight/llgamma/all_conv/plots/kinematics_eta.pdf}
  \caption{$\abseta$ (events)}
\end{subfigure}\hfill
\begin{subfigure}[b]{0.48\textwidth}
  \includegraphics[width=\textwidth]{iso_tight/llgamma/all_conv/plots/kinematics_eta_norm.pdf}
  \caption{$\abseta$ (normalised)}
\end{subfigure}
\caption{All photons, tight isolation.}
\label{fig:tight-all}
\end{figure}
\clearpage

% ====================================================================
\section{Discussion}
\label{sec:discussion}
% ====================================================================

\subsection{Shape Agreement in $\pT$}

Across all scenarios and isolation working points, the normalised $\pT$
distributions show excellent shape agreement between data and MC.  The ratio
panels are consistent with unity within statistical uncertainties over the
full $\pT$ range.  This confirms that the $\pT$ spectrum of the selected
photons is well modelled by the Sherpa simulation and that the event weight
calculation (after the EVNT-level fix) is correct up to the overall
normalisation factor discussed in Section~\ref{sec:residual}.

\subsection{Photon Conversion Rate Discrepancy}
\label{sec:conversion}

A striking observation is the significantly better data--MC agreement for
\textbf{converted} photons (Figures~\ref{fig:loose-conv},
\ref{fig:tight-conv}) compared to \textbf{unconverted} photons
(Figures~\ref{fig:loose-uc}, \ref{fig:tight-uc}), particularly in the
$\abseta$ distributions.

This is understood as a consequence of a known \textbf{TRT barrel tracker
problem in Run~3}.  The Transition Radiation Tracker (TRT), located in the
barrel region ($|\eta| \lesssim 1.0$), plays a key role in classifying
photons as converted or unconverted.  In Run~3, a malfunction in the TRT
barrel causes:
\begin{itemize}
  \item A \textbf{higher fraction of photons classified as unconverted in
    data} than in MC, because the TRT is less efficient at reconstructing
    conversion tracks.
  \item Conversely, a \textbf{lower rate of converted photons in data} compared
    to MC.
\end{itemize}

A correction was previously applied that reclassifies TRT-converted photons
as unconverted, but results from 2024 data show that this correction is
\textbf{insufficient}.  The residual mis-classification leads to:
\begin{itemize}
  \item An \emph{excess} of unconverted photons in data relative to MC,
    appearing as a normalisation deficit in the unconverted MC prediction.
  \item Better agreement for \emph{converted} photons, because this category
    is less affected by the TRT barrel issue (endcap conversions are
    reconstructed differently).
\end{itemize}

This also explains why the $\abseta$ shape discrepancy is concentrated
in the barrel region ($|\eta| < 1.0$), where the TRT barrel is active.

\subsection{Effect of Tight Isolation}

Applying tight isolation instead of loose isolation slightly improves the
overall normalisation agreement, as expected from the rejection of
$Z{+}\text{jets}$ background events in data.  The shape discrepancy in
$\abseta$ for unconverted photons persists, however, confirming that it
is not a background effect but rather the conversion-rate issue described
above.

\subsection{Scale Factors}

The plots include an ``MC (no SFs)'' curve (blue dashed) in addition to
the ``MC (with SFs)'' curve (red).  The scale factors applied are:
lepton reconstruction and identification SFs ($\mathrm{SF}_{\ell_1}$,
$\mathrm{SF}_{\ell_2}$) and photon isolation SF ($\mathrm{SF}_{\gamma,\mathrm{iso}}$).
The difference between the two MC curves is small, indicating that the
scale factors have a modest impact on the kinematic distributions.

% ====================================================================
\section{Summary and Conclusions}
\label{sec:conclusions}
% ====================================================================

\begin{enumerate}
  \item \textbf{Normalisation bug identified and fixed}:
    The DAOD-level sum-of-weights used by the NTupleMaker framework only
    accounts for ${\sim}36\%$ of the true EVNT-level total, inflating MC
    weights by ${\sim}2.8\times$.  Replacing with AMI EVNT-level values
    resolves the over-prediction.

  \item \textbf{Residual ${\sim}20\%$ deficit understood}:
    The remaining MC under-prediction is attributed to missing
    $Z{+}\text{jets}$ background in MC, incomplete MC sample coverage
    (97.59\%), and residual efficiency effects.  This does not affect
    the shower shape analysis, which uses normalised distributions.

  \item \textbf{$\pT$ shapes in good agreement}:
    The normalised $\pT$ distributions show excellent data--MC agreement
    across all scenarios, confirming the MC modelling is adequate for the
    shower shape correction.

  \item \textbf{$\abseta$ discrepancy linked to TRT barrel}:
    The unconverted photon $\abseta$ distribution shows a deficit in MC
    in the barrel region, while converted photons agree much better. This is
    attributed to a known TRT barrel tracker problem in Run~3 that alters
    the photon conversion classification rates between data and MC.  A
    previous correction (reclassifying TRT-converted as unconverted) is
    insufficient for 2024 data.

  \item \textbf{Tight isolation improves agreement marginally}:
    Tightening the isolation requirement reduces background contamination
    and slightly improves the normalisation, but does not resolve the
    conversion-rate discrepancy.
\end{enumerate}

Overall, the kinematic distributions validate the use of the $\Zllg$
sample for the cell-energy reweighting study.  The $\pT$ shape agreement
ensures that the shower shape corrections are not biased by kinematic
mis-modelling.  The conversion-rate discrepancy motivates performing the
shower shape analysis separately for unconverted and converted photons,
and careful attention to the $\abseta$-binned corrections.

\end{document}
